%%%%%%%%%%%%%%%%%%%%%%%%%%%%%%%%%%%%%%%%%
% Homework 2: R / Stata Practice
%%%%%%%%%%%%%%%%%%%%%%%%%%%%%%%%%%%%%%%%%

\documentclass[paper=a4, fontsize=11pt]{scrartcl}

\usepackage[T1]{fontenc}
\usepackage[english]{babel}
\usepackage{amsmath,amsfonts,amssymb}
\usepackage{enumerate}

\usepackage{sectsty}
\allsectionsfont{\normalfont\scshape}

\usepackage[hidelinks]{hyperref}
\hypersetup{colorlinks=true,citecolor=blue,pdfborder=000, urlcolor=red}

\usepackage{fancyhdr}
\pagestyle{fancyplain}
\fancyhead{}
\fancyfoot[L]{}
\fancyfoot[C]{}
\fancyfoot[R]{\thepage}
\renewcommand{\headrulewidth}{0pt}
\renewcommand{\footrulewidth}{0pt}
\setlength{\headheight}{13.6pt}

\setlength\parindent{0pt}

% Horizontal rule
\newcommand{\horrule}[1]{\rule{\linewidth}{#1}}

% Hint box: indented block with label
\newcommand{\hint}[2]{%
  \vspace{0.15cm}%
  \hspace{1em}\textit{\textbf{Hint (#1):}} #2%
  \vspace{0.1cm}%
}

\title{
\normalfont \normalsize
\horrule{0.5pt} \\[0.4cm]
\huge Homework 2: R / Stata Practice \\
\horrule{2pt} \\[0.5cm]
}

\author{Prof. Tzu-Ting Yang}
\date{\normalsize\today}

\begin{document}
\maketitle



%----------------------------------------------------------------------------------------
%	INSTRUCTION
%----------------------------------------------------------------------------------------

\section*{Instruction}

\begin{itemize}
\item[\textbullet] You should submit Homework 2 before 5/27 11:00pm.
\item[\textbullet] There is no late submission (zero points after the deadline).
\item[\textbullet] Homework 2 accounts for 15 points in your final grade.
\item[\textbullet] This homework will help you make progress on your term paper.
\item[\textbullet] You may use \textbf{either Stata or R}. Submit \textbf{two files}:
  \begin{enumerate}
    \item \textbf{Empirical analysis record (HTML)}: A single HTML file documenting your current empirical workflow, generated from your code using markdown.
    \begin{itemize}
      \item \textbf{Stata}: Write code and narrative in a \texttt{.stmd} file and compile with \texttt{markstat using ``filename'', bundle}.
      \item \textbf{R}: Write code and narrative in an \texttt{.Rmd} file and knit to HTML via ``Knit to HTML'' in RStudio or \texttt{rmarkdown::render()}.
    \end{itemize}
    \item \textbf{Answer sheet (PDF)}: A PDF answering all questions (1--6), including tables and figures.
  \end{enumerate}
\item[\textbullet] Recommended R packages: \texttt{haven}, \texttt{dplyr}, \texttt{ggplot2}, \texttt{fixest}, \texttt{rmarkdown}.
\item[\textbullet] Upload files here: \textbf{\url{https://www.dropbox.com/request/1qkUrEkv02ygFaIJcrrX}}
\item[\textbullet] File name format: \texttt{StudentID\_YourName\_HW2.html} and \texttt{StudentID\_YourName\_HW2.pdf}.
\end{itemize}

\bigskip\horrule{0.4pt}\bigskip

%----------------------------------------------------------------------------------------
%	QUESTION 1
%----------------------------------------------------------------------------------------

\subsection*{Question 1 \quad Empirical Analysis Record}

Submit an HTML file that documents your \textbf{current R or Stata code} and results. Generate the HTML from your code using markdown. The HTML should include:
\begin{itemize}
  \item All code you have written so far (data loading, cleaning, visualization, regression).
  \item Output from each code block (tables, figures, regression results).
  \item Brief narrative text explaining what each step does.
\end{itemize}

\medskip
\textit{\textbf{Hint (Stata):}} Create a \texttt{.stmd} file with markdown text and Stata code blocks, then compile:
\begin{quote}\texttt{markstat using ``your\_file'', bundle}\end{quote}
The \texttt{bundle} option embeds all figures directly into the HTML. \\

\textit{\textbf{Hint (R):}} Create an \texttt{.Rmd} file with YAML header \texttt{output: html\_document} and embed R code in \texttt{\`{}\`{}\`{}\{r\}} chunks. Render with:
\begin{quote}\texttt{rmarkdown::render("your\_file.Rmd")}\end{quote}

\bigskip\horrule{0.4pt}\bigskip

%----------------------------------------------------------------------------------------
%	QUESTION 2
%----------------------------------------------------------------------------------------

\subsection*{Question 2 \quad Introduction}

Write a paragraph to introduce your research topic. Include:
\begin{itemize}
  \item A brief description of your research question.
  \item Why you chose to study this topic.
\end{itemize}

\bigskip\horrule{0.4pt}\bigskip

%----------------------------------------------------------------------------------------
%	QUESTION 3
%----------------------------------------------------------------------------------------

\subsection*{Question 3 \quad Sample Construction}

Write a paragraph describing your sample construction process. Include:
\begin{itemize}
  \item Data cleaning steps.
  \item Sample selection criteria.
  \item Final sample size and time period.
\end{itemize}

\bigskip\horrule{0.4pt}\bigskip

%----------------------------------------------------------------------------------------
%	QUESTION 4
%----------------------------------------------------------------------------------------

\subsection*{Question 4 \quad Descriptive Statistics}

Create a table to display descriptive statistics of your sample and provide a brief explanation. Specify this as Table 1 in your answer sheet. Include:
\begin{itemize}
  \item Means and standard deviations for your outcome and control variables.
  \item Comparison between treatment and control groups if applicable.
\end{itemize}

\bigskip\horrule{0.4pt}\bigskip

%----------------------------------------------------------------------------------------
%	QUESTION 5
%----------------------------------------------------------------------------------------

\subsection*{Question 5 \quad Empirical Methodology}

Describe your empirical methodology for estimating causal relationships. Include:
\begin{itemize}
  \item Description of your identification strategy.
  \item Mathematical expressions of your empirical model.
  \item Key assumptions for this method.
\end{itemize}

\bigskip\horrule{0.4pt}\bigskip

%----------------------------------------------------------------------------------------
%	QUESTION 6
%----------------------------------------------------------------------------------------

\subsection*{Question 6 \quad Preliminary Results or AI Assistance}

\textbf{If you have preliminary results}, create a table or figure to present your findings and write a paragraph explaining them. Specify this as Figure 1 (or Table 2) in your answer sheet. Include:
\begin{itemize}
  \item Interpretation of your main estimates.
\end{itemize}

\medskip
\textbf{If you do not have preliminary results yet}, describe how you used generative AI (e.g., ChatGPT, Claude) to assist your research. Include:
\begin{itemize}
  \item What kind of tasks generative AI helped you complete; provide a specific example.
  \item The prompts you gave to the AI.
\end{itemize}

\end{document}
